% ECONOMICS_PROBLEM: {"id": "E42-224", "title": "CBDCs and bank funding", "jel_code": "E42", "source_id": "OP-0329"}
% !TeX program = xelatex
\documentclass[11pt,a4paper]{article}
\usepackage[margin=23mm]{geometry}
\usepackage{fontspec}
\setmainfont{Latin Modern Roman}
\usepackage{amsmath,amssymb,mathtools}
\usepackage{xcolor,enumitem,booktabs,longtable,array,xurl,hyperref}
\definecolor{muted}{HTML}{53636C}
\hypersetup{colorlinks=true,urlcolor=blue,linkcolor=blue}
\setlength{\parindent}{0pt}
\setlength{\parskip}{5pt}
\newcommand{\R}{\mathbb R}
\newcommand{\E}{\mathbb E}
\newcommand{\N}{\mathbb N}
\newcommand{\ind}{\mathbf 1}
\newcommand{\Var}{\operatorname{Var}}
\newcommand{\Cov}{\operatorname{Cov}}
\newcommand{\supp}{\operatorname{supp}}
\DeclareMathOperator*{\argmin}{arg\,min}
\DeclareMathOperator*{\argmax}{arg\,max}
\newcommand{\entrymeta}[3]{{\sffamily\small\color{muted}\textbf{\texttt{#1}}\quad|\quad #2\par #3\par}}
\newcommand{\Problemref}[1]{\texttt{#1}}
\newcommand{\JELref}[2]{\texttt{#2} (JEL \texttt{#1})}
\begin{document}
\section*{CBDCs and bank funding}
\textbf{Problem ID:} \texttt{E42-224}\quad \textbf{JEL:} \texttt{E42}\par
% BEGIN ECONOMICS STATEMENT
\label{problem:OP-0329}
{\sffamily\small\color{muted}Original subject group: Banking, credit and financial stability\par}
\label{definitions:problem:30007611}
\textbf{Audit classification:} Published research agenda. The agenda supports CBDC design and wider effects. The particular deposit/credit/run intervention is editorial; no one-direction general bank-disintermediation theorem is inferred.
\subsection*{Definitions and precise research target}
Fix households choosing among bank deposits, cash and a retail central-bank digital currency. Banks make loans financed by deposits, wholesale funding and equity. A digital-currency design \(a=(\bar q,r,\kappa)\) specifies holding cap \(\bar q\), remuneration rule \(r\) and conversion friction \(\kappa\). Households respond to payment convenience, return, privacy and perceived default risk; banks adjust deposit pricing and lending. Let \(D(a,x)\), \(L(a,x)\) and \(R(a,x)\) be equilibrium deposits, productive credit and a declared crisis-run statistic in state \(x\). The task is to identify or bound the design-dependent changes
\[(D(a,x)-D(0,x),\ L(a,x)-L(0,x),\ R(a,x)-R(0,x)).\]
The baseline \(0\) is the otherwise identical economy without retail digital currency. Distinguish slow portfolio substitution in ordinary states from rapid crisis conversion. Specify whether the central bank replaces lost bank funding and on what collateral terms, since this changes credit effects. Estimate adoption and conversion behavior using justified preference and information assumptions, and characterize welfare-maximizing caps or rates over a feasible set. A result under one bank-funding replacement assumption or one calibrated holding cap does not establish a general disintermediation effect.

\textbf{Additional formal definitions and scope (editorial).} A retail central-bank digital currency is a household-accessible digital claim on the central bank denominated in the domestic unit of account. A holding cap \(\bar q\geq0\) applies per specified person/account aggregation; remuneration \(r\) may depend on holdings and date, and conversion friction \(\kappa\) is a declared fee, delay or limit. Baseline \(a=0\) means no CBDC, not the numerical zero vector of an existing design. Define productive credit \(L\) by a predeclared loan classification and run statistic, for example \(R(a,x)=\Pr(T^{\rm withdrawal}(a)\leq H\mid x)\). Deposits and loans are measured in the same price units. Central-bank replacement funding has explicit collateral, haircut, price and eligibility rules; consolidated fiscal exposure and cash substitution must be included in the comparison.

For the identification part, declare an admissible class \(\Theta\) of structural models, including preferences, technologies, shock laws, measurement/sampling rules and any equilibrium selector. If \(O\) denotes the complete observable history and \(T(\theta)\) the target defined above, the identified set is \(\mathcal I_T=\{T(\theta):\theta\in\Theta,\ \mathcal L_\theta(O)=\mathcal L_{\rm obs}(O)\}\); point identification means this set is a singleton. A \(do(a)\) comparison replaces stated policy/technology equations while fixing the declared exogenous law and re-solving the remaining equations. These definitions do not assert that an identification theorem holds without further restrictions.
\subsection*{Variants and assumption changes}
\begin{enumerate}
\item \textbf{No funding replacement (editorial).} A CBDC competes with deposits while the central bank provides no compensating bank credit; identify deposit and loan changes.
\item \textbf{Collateralized replacement and stress conversion (editorial).} Specify central-bank funding terms and rapid conversion during stress; optimize caps/remuneration jointly with collateral and eligibility.
\end{enumerate}
\subsection*{References and source audit}
\begin{enumerate}
\item Official January 9, 2026 web version; locator: Financial System / 2026 digital-assets priority; Central banking toolkit. Broad agenda; exact model editorial. URL: \url{https://www.bankofengland.co.uk/research/bank-of-england-agenda-for-research}.
\end{enumerate}
\begin{enumerate}\raggedright
\item\hypertarget{definitions:source:30007611:1}{} Bank of England. 2026. \emph{Bank of England Agenda for Research, 2025–2028}. Financial System / 2026 digital-assets priority; Central banking toolkit.
\url{https://www.bankofengland.co.uk/research/bank-of-england-agenda-for-research}.
\end{enumerate}

\subsection*{Common conventions: Source mathematical conventions}

These conventions apply to each entry unless explicitly replaced.

\textbf{Models and identification.} A primitive model \(m\) specifies all probability laws, feasible choices, information, equilibrium and measurement rules used by its target. The admissible class \(\mathcal M\) is fixed independently of outcomes. If \(P_O(m)\) is its observable probability law and \(\tau(m)\) the target, its identified set at observable law \(P\) is
\[
 \mathcal I_\tau(P)=\{\tau(m):m\in\mathcal M,\ P_O(m)=P\}.
\]
Point identification means that this set is a singleton. An empty set signals incompatibility of data and assumptions. Coordinate bounds need not describe the joint set. Bounds are sharp only if every retained target is attainable or arbitrarily approachable by compatible models. Finite bounds require explicit restrictions on unobserved quantities; unrestricted confounding, hidden wealth, extrapolation or arbitrary equilibrium selection can yield uninformative sets.

\textbf{Causal interventions.} A potential outcome \(Y_i(p)\) is the outcome under a complete policy regime \(p\), including timing, financing, information and market spillovers. The target population and units are fixed before treatment. With interference, use the complete assignment vector or a justified exposure mapping; individual no-interference notation alone cannot identify an economy-wide intervention. A difference of mean potential outcomes does not require the cross-world joint distribution. Conditioning on post-treatment migration, survival, enrollment or employment changes the estimand and needs separate justification.

\textbf{Statistical statements.} An estimator, confidence set, forecast or test specifies a sampling law, model class and loss/coverage criterion. Uniform \((1-\alpha)\)-coverage means \(\inf_{m\in\mathcal M}\Pr_m\{\tau(m)\in C_n\}\geq1-\alpha\), or an explicitly asymptotic version. The whole parameter space trivially has coverage, so informativeness requires a stated width, risk or power objective. A forecasting improvement is relative to a named benchmark, information set and evaluation population.

\textbf{Equilibrium and optimization.} An equilibrium comprises feasible plans, optimality, beliefs where needed, budget constraints and all market/resource clearing. The primitives or a stated benchmark define each correspondence \(\mathcal E(p;m)\). Nonemptiness is assumed or proved, never inferred just from compact policy sets. A selected-equilibrium target uses a declared selection rule; a robust target quantifies over all permitted equilibria. Use suprema or infima unless attainment is established. An \(\varepsilon\)-optimal policy is feasible and within \(\varepsilon\) of the finite optimum value. Report an empty feasible set separately.

\textbf{Welfare and accounting.} Normative utility, interpersonal normalization, social weights, numeraire, discounting and the use of public receipts are primitives. Behavioral utility need not coincide with the welfare criterion. Household welfare includes ownership income and rents once; adding profits again to compensating variation generally double-counts them. A monetary equivalent is defined by an explicit transfer and price convention, with monotonicity and range conditions. Average effects, mechanism contributions, transfers and welfare losses are different targets. Infinite sums require convergence and dynamic feasible plans require terminal or no-Ponzi conditions.

\textbf{Source versions.} A working-paper date, revision date and journal publication year are distinguished. A DOI for a journal version is not automatically the DOI of an earlier manuscript. Locators refer to the stated version and printed pages where specified. A metadata-only check does not verify a claim about a full-text conclusion. Every entry's source subsection narrows the provenance claim accordingly.
% END ECONOMICS STATEMENT
\end{document}
